% ════════════════════════════════════════════
% 模型假设（3.模型假设.tex）写作规则 —— 模板内嵌，AI 先读本文件再写
% ════════════════════════════════════════════
% 【结构】两段式：第一段引文 + \begin{itemize} 分点列出每条假设。
% 【引文】第一段固定：「为简化问题，本文做出以下假设：」。
% 【格式】每条假设以 \textbf{假设N：} 开头（加粗+冒号），用分点符号（itemize）区分。
% 【数量】假设条数 = 问题数量（每问对应一条核心假设）。
% 【长度】每条 \textbf{假设N：} ≤2 行，一句话说清该假设的适用前提与理由。
% ════════════════════════════════════════════
\section{模型假设}

为简化问题，本文做出以下假设：

\begin{itemize}[itemindent=2em]
\item \textbf{假设1：} 附件 A 各训练配方的 17 个领域配比与各验证域交叉熵损失之间存在线性混合关系，且该关系在不同参数规模间可由同一组系数近似刻画（配比—损失线性外推假设）；该假设的依据是 1M 训练集平均 $R^2=0.629$ 且 60M/1B 检验集 Pearson 相关达 $0.782/0.651$。
\item \textbf{假设2：} 附件 B 的 $Q_{\text{score}}$ 字段即为数据质量（越大越好，$Q=1$ 为最优），且验证损失随参数量、数据量与质量的变化可用形如 $L=E+AN^{-\alpha}+BD^{-\beta}+C(1-Q)^{\gamma}$ 的广义标度律描述（质量缺口可加假设）；该方向由控制 $N,D$ 后 B6/B7 的组内相关 $-0.925/-0.918$ 一致确认，该形式在四类候选中拟合优度与信息准则最优。
\item \textbf{假设3：} 训练总成本严格按赛题公式分解为基础训练 $6ND$、质量提升 $D[g(Q)-g(Q_0)]_+$ 与长文本注意力 $\eta NDL_{\text{ctx}}$ 三部分之和，其中 $g(Q)$ 取附录 B 规定的三型之一、$\eta=2\times10^{-4}$、$L_{\text{ctx}}$ 由附件 C7 外生给定（成本可加分解假设）；该假设不做额外简化，$C_Q$ 与训练 token 数成正比。
\item \textbf{假设4：} 开源大语言模型能力前沿随时间先加速上升后趋于饱和，可用 logistic 饱和函数刻画；同时规模扩张与非规模技术进步对能力得分的贡献可由面板回归线性分离（技术演进可分解假设）；该假设的先验依据是六维得分存在上界 100，且年份与参数量在控制年份效应后仍可分别识别。
\end{itemize}
